\documentclass[../../main.tex]{subfiles}
\begin{document}
\section{拓扑空间}\label{section:appA-topological-spaces}

我们从定义开始。设 $X$ 是一个集合。$X$ 上的一个\emph{拓扑}是 $X$ 的子集的一个族 $\mathcal T$，其中的子集称为\emph{开子集}，满足
\begin{enumerate}[(i)]
\item $X$ 与 $\varnothing$ 是开集。
\item 任意一族开子集的并是开集。
\item 任意有限族开子集的交是开集。
\end{enumerate}
由集合 $X$ 连同 $X$ 上一个拓扑 $\mathcal T$ 组成的对 $(X,\mathcal T)$ 称为一个\emph{拓扑空间}。通常，当拓扑已清楚时，我们略去对它的提及，简说“$X$ 是一个拓扑空间”。

与拓扑空间相关联的构造和定义有很多。这里我们概述本书中最重要的那些。

设 $X$ 是拓扑空间，$p\in X$，$S\subset X$。
\begin{itemize}
\item $p$ 的一个\emph{邻域}是包含 $p$ 的开子集。类似地，集合 $S$ 的一个邻域是包含 $S$ 的开子集。（注意：有些作者在更一般的意义下使用“邻域”一词，即指包含某个含 $p$ 或 $S$ 的开子集的子集。）
\item 如果 $X\setminus S$ 是开集（其中 $X\setminus S$ 表示差集 $\{x\in X:x\notin S\}$），则称 $S$ 是\emph{闭的}。
\item $S$ 的\emph{内部}记作 $\operatorname{Int}S$，是 $X$ 中包含于 $S$ 的所有开子集之并。
\item $S$ 的\emph{外部}记作 $\operatorname{Ext}S$，是 $X$ 中包含于 $X\setminus S$ 的所有开子集之并。
\item $S$ 的\emph{闭包}记作 $\bar S$，是 $X$ 中包含 $S$ 的所有闭子集之交。
\item $S$ 的\emph{边界}记作 $\partial S$，是 $X$ 中既不属于 $\operatorname{Int}S$ 也不属于 $\operatorname{Ext}S$ 的所有点的集合。
\item 如果 $p\in S$ 且 $p$ 有邻域 $U\subset X$ 使 $U\cap S=\{p\}$，则称 $p$ 是 $S$ 的一个\emph{孤立点}。
\item 如果 $p\in X$（不必在 $S$ 中）的每个邻域都至少含有 $S$ 中一个不同于 $p$ 的点，则称 $p$ 是 $S$ 的一个\emph{极限点}。
\item 如果 $\bar S=X$，等价地，如果 $X$ 的每个非空开子集都至少含有 $S$ 的一个点，则称 $S$ 在 $X$ 中\emph{稠密}。
\item 如果 $\bar S$ 不含任何非空开子集，则称 $S$ 在 $X$ 中\emph{无处稠密}。
\end{itemize}

拓扑学中最重要的概念是连续映射与收敛序列，我们下面给出定义。设 $X$ 与 $Y$ 是拓扑空间。
\begin{itemize}
\item 如果对每个开子集 $U\subset Y$，原像 $F^{-1}(U)$ 在 $X$ 中开，则称映射 $F:X\to Y$ \emph{连续}。
\item 连续双射且逆连续的映射 $F:X\to Y$ 称为\emph{同胚}。如果存在从 $X$ 到 $Y$ 的同胚，我们就说 $X$ 与 $Y$ \emph{同胚}。
\item 如果连续映射 $F:X\to Y$ 的每个点 $p\in X$ 都有邻域 $U\subset X$，使 $F(U)$ 在 $Y$ 中开且 $F$ 限制为从 $U$ 到 $F(U)$ 的同胚，则称 $F$ 是\emph{局部同胚}。
\item 给定 $X$ 中点的序列 $(p_i)_{i=1}^\infty$ 和点 $p\in X$，如果对 $p$ 的每个邻域 $U$，都存在正整数 $N$ 使对所有 $i\geqslant N$ 都有 $p_i\in U$，则称该序列\emph{收敛到} $p$。在此情形中，写作 $p_i\to p$ 或 $\lim_{i\to\infty}p_i=p$。
\end{itemize}

\begin{exercise}\label{exercise:appA-A-1}
设 $F:X\to Y$ 是拓扑空间之间的映射。证明下列每条性质都等价于 $F$ 的连续性：
\begin{enumerate}[(a)]
\item 对每个子集 $A\subset X$，$F(\bar A)\subset\overline{F(A)}$。
\item 对每个子集 $B\subset Y$，$F^{-1}(\operatorname{Int}B)\subset\operatorname{Int}F^{-1}(B)$。
\end{enumerate}
\end{exercise}

\begin{exercise}\label{exercise:appA-A-2}
设 $X,Y,Z$ 是拓扑空间。证明下列映射连续：
\begin{enumerate}[(a)]
\item 恒等映射 $\operatorname{Id}_X:X\to X$，定义为 $\operatorname{Id}_X(x)=x$ 对所有 $x\in X$。
\item 任意常值映射 $F:X\to Y$（即满足 $F(x)=F(y)$ 对所有 $x,y\in X$ 的映射）。
\item 连续映射 $F:X\to Y$ 与 $G:Y\to Z$ 的任意复合 $G\circ F$。
\end{enumerate}
\end{exercise}

\begin{exercise}\label{exercise:appA-A-3}
设 $X$ 与 $Y$ 是拓扑空间。假设 $F:X\to Y$ 连续且在 $X$ 中 $p_i\to p$。证明在 $Y$ 中 $F(p_i)\to F(p)$。
\end{exercise}

拓扑空间最重要的例子——我们的绝大多数流形例子都以某种方式由它们构造——描述如下。

\begin{example}[离散空间]\label{example:appA-discrete-spaces}
如果 $X$ 是任意集合，$X$ 上的\emph{离散拓扑}是这样定义的拓扑：宣告 $X$ 的每个子集都是开集。凡带有离散拓扑的空间都称为\emph{离散空间}。
\end{example}

\begin{example}[度量空间]\label{example:appA-metric-spaces}
一个\emph{度量空间}是一个集合 $M$ 连同一个距离函数（也叫度量）$d:M\times M\to\mathbb R$（其中 $\mathbb R$ 表示实数集），它对所有 $x,y,z\in M$ 满足下列性质：
\begin{enumerate}[(i)]
\item \textbf{正性：}$d(x,y)\geqslant0$，等号当且仅当 $x=y$ 时成立。
\item \textbf{对称性：}$d(x,y)=d(y,x)$。
\item \textbf{三角不等式：}$d(x,z)\leqslant d(x,y)+d(y,z)$。
\end{enumerate}
如果 $M$ 是度量空间，$x\in M$，$r>0$，围绕 $x$ 半径为 $r$ 的\emph{开球}是集合
\begin{align*}
B_r(x)=\{y\in M:d(x,y)<r\},
\end{align*}
半径为 $r$ 的\emph{闭球}是
\begin{align*}
\bar B_r(x)=\{y\in M:d(x,y)\leqslant r\}.
\end{align*}
$M$ 上的\emph{度量拓扑}定义为：宣告子集 $S\subset M$ 是开的，如果对每个点 $x\in S$，都存在某个 $r>0$ 使 $B_r(x)\subset S$。
\end{example}

如果 $M$ 是度量空间而 $S$ 是 $M$ 的任意子集，距离函数限制到 $S$ 中点对，就把 $S$ 变成一个度量空间，从而也是一个拓扑空间。我们对度量空间使用下列标准术语：
\begin{itemize}
\item 如果存在正数 $R$ 使对所有 $x,y\in S$ 都有 $d(x,y)\leqslant R$，则称子集 $S\subset M$ \emph{有界}。
\item 如果 $S$ 是 $M$ 的非空有界子集，$S$ 的\emph{直径}是数
\begin{align*}
\operatorname{diam}S=\sup\{d(x,y):x,y\in S\}.
\end{align*}
\item 如果对每个 $\varepsilon>0$，都存在整数 $N$ 使 $i,j\geqslant N$ 蕴含 $d(x_i,x_j)<\varepsilon$，则称 $M$ 中点的序列 $(x_i)_{i=1}^\infty$ 是一个\emph{Cauchy 序列}。
\item 如果 $M$ 中每个 Cauchy 序列都收敛到 $M$ 中一点，则称度量空间 $M$ \emph{完备}。
\end{itemize}

\begin{example}[Euclid 空间]\label{example:appA-euclidean-spaces}
对每个整数 $n\geqslant1$，实 $n$ 元组的集合 $\mathbb R^n$ 称为 $n$ 维\emph{Euclid 空间}。我们用 $(x^1,\dots,x^n)$、$(x^i)$ 或 $x$ 表示 $\mathbb R^n$ 中的一点；数 $x^i$ 称为 $x$ 的\emph{分量}或\emph{坐标}。（当 $n$ 较小时，我们常使用更传统的坐标名称，如 $(x,y,z)$。）注意我们用上标而非下标写点 $(x^1,\dots,x^n)\in\mathbb R^n$ 的坐标——线性代数和微积分教材通常用下标——这是为了与第 1 章解释的 Einstein 求和约定一致。按约定，$\mathbb R^0$ 是单元素集 $\{0\}$。

对每个 $x\in\mathbb R^n$，$x$ 的\emph{Euclid 范数}是非负实数
\begin{align*}
|x|=\sqrt{(x^1)^2+\cdots+(x^n)^2},
\end{align*}
而对 $x,y\in\mathbb R^n$，\emph{Euclid 距离函数}定义为
\begin{align*}
d(x,y)=|x-y|.
\end{align*}
这个距离函数把 $\mathbb R^n$ 变成一个完备度量空间。由此得到的 $\mathbb R^n$ 上的度量拓扑称为\emph{Euclid 拓扑}。
\end{example}

\begin{example}[复 Euclid 空间]\label{example:appA-complex-euclidean-spaces}
我们有时也有机会处理复 Euclid 空间。我们把复数集 $\mathbb C$ 作为集合简单地看作 $\mathbb R^2$，复数 $x+\mathrm{i}y$ 对应于 $(x,y)\in\mathbb R^2$。对任意正整数 $n$，$n$ 维\emph{复 Euclid 空间}是复 $n$ 元组的集合 $\mathbb C^n$。当通过对应
\begin{align*}
(x^1+\mathrm{i}y^1,\dots,x^n+\mathrm{i}y^n)\leftrightarrow(x^1,y^1,\dots,x^n,y^n)
\end{align*}
等同于 $\mathbb R^{2n}$ 时，它就成为一个拓扑空间。
\end{example}

\begin{example}[Euclid 空间的子集]\label{example:appA-subsets-of-euclidean}
$\mathbb R^n$ 或 $\mathbb C^n$ 的每个子集在赋予 Euclid 度量后都成为一个度量空间，从而也是一个拓扑空间。每当我们提到这样一个子集时，除非另行说明，总假定它带此度量拓扑。它是完备度量空间，当且仅当它是 $\mathbb R^n$ 的闭子集。下面是我们经常处理的 Euclid 空间的一些标准子集：
\begin{itemize}
\item \emph{单位区间}是 $\mathbb R$ 的子集 $I$，定义为
\begin{align*}
I=[0,1]=\{x\in\mathbb R:0\leqslant x\leqslant1\}.
\end{align*}
\item $n$ 维\emph{（开）单位球}是 $\mathbb R^n$ 的子集 $B^n$，定义为
\begin{align*}
B^n=\{x\in\mathbb R^n:|x|<1\}.
\end{align*}
\item $n$ 维\emph{闭单位球}是 $\mathbb R^n$ 的子集 $\bar B^n$，定义为
\begin{align*}
\bar B^n=\{x\in\mathbb R^n:|x|\leqslant1\}.
\end{align*}
术语（开）单位圆盘和闭单位圆盘通常分别用于 $B^2$ 和 $\bar B^2$。
\item 对 $n\geqslant0$，\emph{（单位）$n$ 球面}是 $\mathbb R^{n+1}$ 的子集 $S^n$，定义为
\begin{align*}
S^n=\{x\in\mathbb R^{n+1}:|x|=1\}.
\end{align*}
有时借助 $\mathbb C^{n+1}$ 与 $\mathbb R^{2n+2}$ 的通常等同，把奇维球面 $S^{2n+1}$ 看作 $\mathbb C^{n+1}$ 的子集是有用的。
\item \emph{（单位）圆}是 $1$-球面 $S^1$，既可看作 $\mathbb R^2$ 的子集，也可看作 $\mathbb C$ 的子集。
\end{itemize}
\end{example}

\subsection*{Hausdorff 空间}

拓扑空间使我们能描述各种各样的“空间”概念。但对流形理论而言，任意拓扑空间过于一般，因为它们可能有一些令人不快的性质，如下一个习题所示。

\begin{exercise}\label{exercise:appA-A-9}
设 $X$ 是任意集合。证明 $\{X,\varnothing\}$ 是 $X$ 上的一个拓扑，称为\emph{平凡拓扑}。证明当 $X$ 赋予此拓扑时，$X$ 中每个序列都收敛到 $X$ 的每个点，且从任意拓扑空间到 $X$ 的每个映射都连续。
\end{exercise}

为避免这类因 $X$ 没有足够多开子集而产生的病态情形，我们常把注意力限制在满足下列特殊条件的拓扑空间上。如果对每对不同的点 $p,q\in X$，都存在不相交的开子集 $U,V\subset X$ 使 $p\in U$ 且 $q\in V$，则称拓扑空间 $X$ 是一个\emph{Hausdorff 空间}。

\begin{exercise}\label{exercise:appA-A-10}
证明每个度量空间在度量拓扑下都是 Hausdorff 的。
\end{exercise}

\begin{exercise}\label{exercise:appA-A-11}
设 $X$ 是 Hausdorff 空间。证明 $X$ 的每个有限子集都是闭的，且 $X$ 中每个收敛序列都有唯一极限。
\end{exercise}

\subsection*{基与可数性}

设 $X$ 是拓扑空间。如果 $X$ 的每个开子集都是 $\mathcal B$ 中某些元素之并，则称 $X$ 的开子集族 $\mathcal B$ 是 $X$ 拓扑的一个\emph{基}。

更一般地，设 $X$ 只是一个集合，$\mathcal B$ 是 $X$ 的子集族，满足下列条件：
\begin{enumerate}[(i)]
\item $X=\bigcup_{B\in\mathcal B}B$。
\item 如果 $B_1,B_2\in\mathcal B$ 且 $x\in B_1\cap B_2$，则存在 $B_3\in\mathcal B$ 使 $x\in B_3\subset B_1\cap B_2$。
\end{enumerate}
于是 $\mathcal B$ 中元素的所有并所成族是 $X$ 上的一个拓扑，称为由 $\mathcal B$ \emph{生成的拓扑}，而 $\mathcal B$ 是此拓扑的一个基。

如果 $X$ 是拓扑空间而 $p\in X$，$p$ 处的一个\emph{邻域基}是 $p$ 的邻域的一个族 $\mathcal B_p$，使 $p$ 的每个邻域都至少含有 $\mathcal B_p$ 中的一个 $B$。

一个集合如果能与正整数集建立双射，就称为\emph{可数无限}；如果是有限或可数无限，就称为\emph{可数}。如果每个点处都有可数邻域基，就称拓扑空间 $X$ \emph{第一可数}；如果其拓扑有可数基，就称 $X$ \emph{第二可数}。由于 $X$ 的一个可数基在每点处都含有一个可数邻域基，第二可数性蕴含第一可数性。

下一个引理表达第一可数空间最重要的性质。说一个序列\emph{终于}某个子集，是指该序列除有限多项外都在该子集中。

\begin{lemma}[序列引理]\label{lemma:appA-sequence-lemma}
设 $X$ 是第一可数空间，$A\subset X$ 是任意子集，$x\in X$。
\begin{enumerate}[(a)]
\item $x\in\bar A$，当且仅当 $x$ 是 $A$ 中点的某个序列的极限。
\item $x\in\operatorname{Int}A$，当且仅当 $X$ 中每个收敛到 $x$ 的序列都终于 $A$。
\item $A$ 在 $X$ 中闭，当且仅当 $A$ 含有 $A$ 中点的每个收敛序列的每个极限。
\item $A$ 在 $X$ 中开，当且仅当 $X$ 中每个收敛到 $A$ 中一点的序列都终于 $A$。
\end{enumerate}
\end{lemma}

\begin{exercise}\label{exercise:appA-A-13}
证明序列引理。
\end{exercise}

\begin{exercise}\label{exercise:appA-A-14}
证明每个度量空间都是第一可数的。
\end{exercise}

\begin{exercise}\label{exercise:appA-A-15}
证明 $\mathbb R^n$ 中半径为有理数、中心坐标为有理数的所有开球所成集合，是 Euclid 拓扑的一个可数基，因而 $\mathbb R^n$ 是第二可数的。
\end{exercise}

第二可数空间最重要的性质之一由下列命题表达。设 $X$ 是拓扑空间。$X$ 的一个\emph{覆盖}是 $X$ 的子集的一个族 $\mathcal U$，其并为 $X$；如果 $\mathcal U$ 中每个集合都开，就称为\emph{开覆盖}。$\mathcal U$ 的一个\emph{子覆盖}是 $\mathcal U$ 的一个仍为覆盖的子族。

\begin{proposition}\label{proposition:appA-second-countable-subcover}
设 $X$ 是第二可数拓扑空间。$X$ 的每个开覆盖都有可数子覆盖。
\end{proposition}

\begin{proof}
设 $\mathcal B$ 是 $X$ 的一个可数基，$\mathcal U$ 是 $X$ 的任意开覆盖。令 $\mathcal B'\subset\mathcal B$ 为满足“存在某个 $U\in\mathcal U$ 使 $B\subset U$”的基开子集 $B\in\mathcal B$ 所成族。对每个 $B\in\mathcal B'$，选取一个含 $B$ 的特定集合 $U_B\in\mathcal U$。族 $\{U_B:B\in\mathcal B'\}$ 可数，故只需证明它覆盖 $X$。给定点 $x\in X$，存在某个 $V\in\mathcal U$ 含 $x$；由于 $\mathcal B$ 是基，存在 $B\in\mathcal B$ 使 $x\in B\subset V$。这特别蕴含 $B\in\mathcal B'$，因而 $x\in B\subset U_B$。
\end{proof}
\end{document}
